\section{Pendahuluan}
\label{sec:intro}

Gelar sarjana di Indonesia diakhiri dengan \emph{skripsi}, karya tulis akhir yang bab-babnya mengikuti konvensi program studi: bab metode (\emph{BAB~III}) yang subbabnya hampir tetap di beberapa program, bab hasil (\emph{BAB~IV}) yang bentuknya mengikuti desain penelitian, statistik yang dihitung di SPSS, Excel atau Python, tinjauan pustaka yang harus mengutip sumber nyata, dan bahasa Indonesia baku yang mengikuti aturan ejaan nasional (EYD) dan kamus (KBBI). Mahasiswa, pembimbing dan alat yang dibuat untuk mereka semakin menyerahkan bagian-bagian pekerjaan ini kepada model bahasa besar (LLM). Model mana yang sebaiknya mereka gunakan?

Pertanyaannya menjadi lebih sulit, bukan lebih mudah. Satu SDK vendor kini menyajikan puluhan model dari beberapa vendor, dengan harga per panggilan yang berbeda hingga dua ordo besaran dan sering kali dengan pengaturan reasoning effort; model open-weight dijalankan pada GPU laboratorium dengan biaya sebesar harga listrik. Papan peringkat publik mengurutkan model berdasarkan pengetahuan dan penalaran bahasa Inggris~\cite{hendrycks2021mmlu,liang2023helm}, dan benchmark Indonesia mengukur pemahaman bahasa, sentimen, dan pengetahuan ujian sekolah~\cite{wilie2020indonlu,koto2020indolem,winata2023nusax,koto2023indommlu,wang2024seaeval}. Tidak satu pun dari mereka menanyakan apakah sebuah model membuat kerangka bab metode sebagaimana diharapkan oleh program studi agribisnis atau program studi hukum keluarga Islam, hanya mengutip sumber yang diberikan padanya, mencantumkan rujukan yang ada, menulis abstrak dalam bahasa Indonesia baku tanpa mengarang angka, atau menghitung uji Mann--Whitney yang benar dari sebuah berkas CSV. Demikian pula mereka tidak melaporkan berapa biaya suatu jawaban atau model mana yang menghasilkan jawaban itu---bukan pertanyaan sepele, karena sebuah SDK vendor dapat diam-diam melayani model yang berbeda dari model yang diminta~\cite{arya2026paperagent}.

Kami menyajikan \emph{AcademicBench}, delapan tugas yang diambil dari alur kerja skripsi, dan menggunakannya untuk membandingkan \NApi{} model API dari satu SDK, yang kami kelompokkan ke dalam dua tier berdasarkan nama, dengan \NLocal{} model lokal menggunakan prompt bahasa Indonesia yang identik. Kontribusi kami adalah:
\begin{itemize}
  \item \textbf{Tugas dan butir} (Bagian~\ref{sec:bench}): \NItems{} butir: 225 butir publik yang baru ditulis dan 30 cuplikan privat dari skripsi nyata (dipublikasikan sebagai hash saja); penskoran deterministik untuk enam tugas dan untuk separuh deterministik dari dua tugas prosa, separuh lainnya berasal dari juri LLM dari vendor lain selain vendor kandidat, dengan kalibrasi kelunakan dan analisis reliabilitas.
  \item \textbf{Pencatatan} (Bagian~\ref{sec:setup}): setiap panggilan mencatat model yang dilaporkan oleh SDK sebagai yang menjawab, reasoning effort, token, dan AI units vendor (dikonversi ke US\$); pelaksanaan lokal mencatat waktu GPU (dikonversi dengan harga on-demand untuk GPU yang sama) dan energi papan GPU yang diukur.
  \item \textbf{Temuan} (Bagian~\ref{sec:results}): tiga model API terbaik tidak berbeda signifikan, sementara biaya per butir model teratas adalah \CostRatio{} kali dari model \CheapGap{} di bawahnya; pertanyaan metodologi dan bahasa Indonesia baku mendekati batas atas untuk model API, sedangkan rujukan nyata, kerangka spesifik program studi, dan perbaikan LaTeX membedakan mereka; sebagian besar rujukan yang tidak ditemukan di Crossref juga tidak ditemukan di basis data lain mana pun; rata-rata, model lokal tertinggal dibanding model API pada setiap tugas, paling jauh pada rujukan dan kode, paling dekat pada bahasa Indonesia baku.
  \item \textbf{Sebuah kit terbuka}: runner dengan gerbang biaya dan batas anggaran, pemberi skor, sandbox, audit rujukan, butir publik, skor per butir dan pembuat lembar penilaian tersamar untuk penilai manusia, dalam repositori paper-agent (\url{https://github.com/devnolife/paper-agent}).
\end{itemize}
